\scanpage{145}
et $x^*\in V(N)(A)=\operatorname{Hom}_{k\text{-mod}}(N,A)$,
\diagram{
\begin{tikzcd}
u^*(x^*)=\widetilde{x}^*\circ u & M\arrow[r,"u"] & \operatorname{Sym}^*N\arrow[r,"\widetilde{x}^*"] & A\\
\operatorname{Hom}_{k\text{-mod}}(M,A)\arrow[u,phantom,"\in" description,sloped] & & &
\end{tikzcd}
}
où $\widetilde{x}^*$ est l'hom de $k$-Alg. \unsure{prolgt} can. de $x^*$.

On a donc
\[
u^*(x^*)(\deleted{\unsure{x}}\inserted{m})=
\sum_{i\geq0}\widetilde{x}^{*i}(u_i(\unsure{x|m}))
\qquad(\widetilde{x}^{*i}\text{ étant la }\deleted{\unsure{exten-}}\inserted{restr.}\text{ de }x^*\text{ à }\operatorname{Sym}^iN)
\]
d'où, si $\lambda\in A$
\[
u^*(\lambda x^*)(m)=\sum_{i\geq0}\lambda^i\widetilde{x}^{*i}(u_i(m))
=\sum_{i\geq0}(\deleted{\unsure{\lambda^i}}\widetilde{x}^{*i})\unsure{((\lambda^iu_i)m)}
\]
Donc dans la correspondance entre $u^*$ et $u$, sur $k$ ou après chgt. de base sur un $A$, la \deleted{multiplication}\inserted{composition} \deleted{\unsure{par $u$}} de $x^*$ avec une homothétie \deleted{$\unsure{\lambda(N_A)=\lambda\in A}$} $\lambda$ dans le $O_k$-module $N$ (ou $O_A$-module $N_A$) correspond à l'opération
\[
u=\sum u_i\longmapsto u^\lambda=\sum\lambda^iu_i.
\]
On en conclut que $u^*$ est \uline{homogène de} degré $n$ (i.e. satisfait, pour $A$, $\lambda$ variables,
\[
:\quad u^*\circ\lambda_{V(N)}=\lambda^n_{V(M)}\circ u^*
\]
ssi tous les $u_i$ pour $i\neq n$ sont nuls, i.e. ssi $u$ est un hom
\[
M\mathbin{\text{---}}\operatorname{Sym}^nN.
\]
[En effet, il faut avoir $\sum(\lambda^i-\lambda^n)u_i=0$ $\forall A,\lambda\in A$, prenant $A=k[T]$ et $\lambda=T$, on doit avoir
\[
\sum_{i\neq n}(T^i-T^n)u_i\unsure{(x)}=0\ ,\quad
\text{i.e. }\sum(T^i-T^n)u_i(x)=0
\]
pour tt $x\in M$, $n\neq i$, d'où \deleted{$u_i(x)$} $u_i(x)=0$
\margin{144}
